\section{The Field at a Glance}\label{sec:landscape}

\paragraph{Growth.} Annual output rose from \nxrPapersSeventeen{} papers in 2017 to
\nxrPapersTwentyFour{} in 2024, a factor of \nxrGrowthFactor{} and a compound growth of \nxrCagrTotal{}
a year (\cref{fig:growth}a). Growth was steady until 2022 and then jumped, largely with the metaverse.
Papers that use the word went from \nxrMetaverseTwentyOne{} in 2021 to \nxrMetaverseTwentyTwo{} in 2022
and \nxrMetaverseTwentyThree{} in 2023. More than half of all included papers (\nxrShareSince{})
appeared from 2023 onward. The field is conference-led: \nxrShareConfPapers{} of papers are conference
papers. It is also spread across communities that seldom publish in each other's venues. Communications
venues (Globecom, ICC, \emph{IEEE Network}, \emph{JSAC}) sit beside multimedia-systems venues (MMSys,
ACM Multimedia, \emph{IEEE Transactions on Multimedia}) and broad open-access journals. The ten largest
venues hold only \nxrShareTopTenVenues{} of the corpus. XR venues proper barely appear, so perceptual
results from the XR community, such as tolerable latency, reach networking papers late.

\begin{figure}[t]
\centering
\includegraphics[width=\linewidth]{fig_growth}
\caption{Growth and evaluation realism. (a)~Included papers per year; 2016, 2025 and 2026 are partial
because of the harvest window and indexing lag. (b)~Share of papers by the most realistic evaluation
they report. Evaluation on real networks or testbeds (the two darkest bands) shrank as simulation
grew.}\label{fig:growth}
\end{figure}

\paragraph{From delivering content to allocating resources.} The composition of the field changed more
than its size (\cref{fig:shifts}). In 2017--19 the largest problem was rate adaptation and streaming,
with \nxrShProbRateAdaptationStreamingEarly{} of papers. By 2022--24 its share had fallen to
\nxrShProbRateAdaptationStreamingLate{}, the largest decline of any problem. Resource allocation rose from
\nxrShProbResourceAllocationEarly{} to \nxrShProbResourceAllocationLate{}, and architecture and vision
papers from \nxrShProbArchitectureVisionEarly{} to \nxrShProbArchitectureVisionLate{}. Security, privacy
and trust grew fastest of all, at \nxrCagrProbSecurityPrivacyTrust{} a year. At the level of functions,
anticipation fell from \fnAnticipateEarly{} to \fnAnticipateLate{} of output and provisioning rose from
\fnProvisionEarly{} to \fnProvisionLate{}. These are shares: streaming still grew in absolute terms, but
far more slowly than the field. The shift reflects a change of authors. The early literature was led by
multimedia-systems groups building 360-degree players. The later literature is led by wireless groups
who formulate XR delivery as a joint allocation of radio, computing and energy, and by authors of 6G and
metaverse visions.

\paragraph{From 360-degree video to ``XR traffic''.} The reference workload changed with the problems.
360-degree video was the content of \nxrShContVideoThreeSixtyEarly{} of papers in 2017--19 and of
\nxrShContVideoThreeSixtyLate{} in 2022--24. The largest gain went to papers that name no specific
content and model XR as a rate and a delay bound (\nxrShContGeneralXrTrafficEarly{} $\to$
\nxrShContGeneralXrTrafficLate{}). Among specific contents, volumetric media, digital-twin data, rendered
frames and holograms grew, while AR overlays and haptic signals declined. The rise of generic traffic
matters for the evidence. A constant-rate flow with a delay bound cannot express viewport dependence,
frame-size variability or the coupling between uplink pose and downlink content. These are exactly the
properties that the content-specific literature has measured.

\begin{figure}[t]
\centering
\includegraphics[width=\linewidth]{fig_shifts}
\caption{Change in composition between 2017--19 (open circles) and 2022--24 (filled circles), with the
difference in percentage points. Colours give the function of \cref{sec:framework}. The field moved from
anticipation (streaming 360-degree video) to provisioning (allocating resources for generic XR
traffic).}\label{fig:shifts}
\end{figure}

\paragraph{From systems to learning, and from testbeds to simulators.} System implementations fell from
\nxrShMethSystemImplementationEarly{} to \nxrShMethSystemImplementationLate{} of papers. Reinforcement
learning grew from \nxrShMethReinforcementLearningEarly{} to \nxrShMethReinforcementLearningLate{}, at
\nxrCagrMethReinforcementLearning{} a year, and supervised deep learning from
\nxrShMethDeepLearningEarly{} to \nxrShMethDeepLearningLate{}. User studies, rare to begin with, fell from
\nxrShMethUserStudyEarly{} to \nxrShMethUserStudyLate{}, and in ten years the corpus contains only
\nxrNDatasets{} papers whose main contribution is a dataset. Evaluation followed the methods
(\cref{fig:growth}b). Evaluation on a real network or testbed fell from \nxrEvalRealEarly{} of papers to
\nxrEvalRealLate{}, while simulation and analysis alone rose from \nxrEvalSimAnaEarly{} to
\nxrEvalSimAnaLate{}. The fastest-growing methods are those least often tested on real links:
\nxrRLSimShare{} of reinforcement-learning papers are evaluated by simulation alone and only
\nxrRLRealShare{} on a real testbed. Learned controllers may well work. But most of their reported gains
are gains inside a model of the network, usually one with synthetic channels, traffic and user motion.

\paragraph{What is new.} Four clusters of author keywords rose fastest between 2018--20 and 2023--25. The
first is the metaverse and 6G (``metaverse'' from \nxrKwMetaverseEarly{} paper to \nxrKwMetaverseLate{};
``6G'' from \nxrKwSixGEarly{} to \nxrKwSixGLate{}). The second is digital twins (\nxrKwDigitalTwinEarly{}
$\to$ \nxrKwDigitalTwinLate{}). The third is learning-based networking, including federated learning
(\nxrKwFederatedLearningEarly{} $\to$ \nxrKwFederatedLearningLate{}). The fourth is new representations
and transmission paradigms: semantic communication (\nxrKwSemanticCommunicationEarly{} $\to$
\nxrKwSemanticCommunicationLate{}), Gaussian splatting (\nxrKwGaussianSplattingEarly{} $\to$
\nxrKwGaussianSplattingLate{}) and holographic communication (\nxrKwHolographicCommunicationEarly{} $\to$
\nxrKwHolographicCommunicationLate{}). Much of this new literature consists of visions and surveys.
Sections~\ref{sec:encode}--\ref{sec:assure} ask how much of it has yet met a real network.
